\section{Resultados}
\label{sec:resultados}

\subsection{Resposta da camada de saída para uma amostra}
\label{sec:amostra}

A Figura~\ref{fig:amostra-saidas} apresenta a resposta completa da
camada de saída para uma única imagem, um dígito 3 do conjunto de
teste, medida na simulação com LM741. O painel esquerdo mostra a
entrada: a imagem $8\times 8$, isto é, as 64 tensões aplicadas às
fontes. O painel direito mostra, para cada um dos dez neurônios de
saída, duas barras: a tensão $V(z_k)$ do somador linear (azul) e a
tensão $V(out_k)$ do comparador correspondente (verde); a linha
tracejada vermelha é o limiar $V_{TH} = \SI{2,65}{\volt}$.

Três observações resumem o funcionamento da rede. Primeiro, apenas
$z_3$ está acima do limiar; as demais saídas estão volts abaixo, a
maioria em valores negativos, resultado do treinamento que empurra os
dígitos rejeitados para baixo. Segundo, os comparadores traduzem essa
separação em níveis lógicos: somente $out_3$ está no nível alto,
próximo de \SI{4,8}{\volt}, e os outros nove permanecem em
\SI{0}{\volt}. Terceiro, a margem entre $z_3$ e o limiar (e entre o
limiar e a segunda maior saída) é da ordem de volts, muito acima dos
erros de milivolts a centenas de milivolts introduzidos pelos
amplificadores reais, o que explica a robustez da decisão observada em
todas as variantes.

\begin{figure}[htb]
  \centering
  \includegraphics[width=\textwidth]{amostra_saidas.png}
  \caption{Resposta da camada de saída para um dígito 3 do conjunto de
  teste, medida no LTspice (variante LM741): tensões $V(z_k)$ dos dez
  somadores e $V(out_k)$ dos dez comparadores, com o limiar $V_{TH}$
  indicado.}
  \label{fig:amostra-saidas}
\end{figure}

\subsection{Sequência de vinte janelas}
\label{sec:sequencia}

A Figura~\ref{fig:ondas} apresenta as formas de onda das dez primeiras
janelas da sequência de teste, na variante LM741. No painel superior
estão as dez saídas $z_k$: os traços cinza são os dez sinais
sobrepostos e, em cada janela, o traço azul destaca a saída do dígito
efetivamente apresentado. O painel inferior mostra a tensão do
comparador desse mesmo dígito, com o rótulo da janela anotado. A cada
troca de imagem observa-se a mesma sequência de eventos: as saídas
$z_k$ acomodam-se nos novos valores em algumas dezenas de
microssegundos (limitadas pelo slew rate do LM741,
Seção~\ref{sec:tempo}), a saída do novo dígito cruza o limiar e o
comparador correspondente comuta para o nível alto, onde permanece até
o fim da janela. O comparador da janela anterior comuta para baixo na
mesma transição.

\begin{figure}[htb]
  \centering
  \includegraphics[width=\textwidth]{formas_de_onda.png}
  \caption{Dez janelas de \SI{250}{\micro\second} da sequência de
  teste (variante LM741). Acima: as dez saídas $z_k$ (cinza), com a
  saída do dígito apresentado em azul e o limiar $V_{TH}$ tracejado.
  Abaixo: o comparador do dígito apresentado em cada janela.}
  \label{fig:ondas}
\end{figure}

A Figura~\ref{fig:heatmap} condensa as vinte janelas em uma única
matriz. Cada coluna é uma janela (o dígito apresentado está no eixo
horizontal) e cada linha é um comparador; a intensidade da célula é a
tensão medida do comparador ao fim da janela, de \SI{0}{\volt} (claro)
a \SI{4,8}{\volt} (escuro). O contorno vermelho marca, em cada coluna,
a célula do dígito correto. O padrão observado é uma diagonal dupla de
células escuras exatamente sobre os contornos (a sequência apresenta
os dígitos 0 a 9 duas vezes) sem nenhuma célula escura fora deles: nas
vinte janelas, o comparador correto disparou sozinho, o que
corresponde a \vDisparosLm~acertos pelo critério de disparo único.

\begin{figure}[htb]
  \centering
  \includegraphics[width=\textwidth]{comparadores_heatmap.png}
  \caption{Tensão dos dez comparadores em cada uma das vinte janelas
  de teste (variante LM741). O contorno vermelho indica o dígito
  apresentado; a diagonal escura sobre os contornos corresponde aos
  vinte disparos únicos corretos.}
  \label{fig:heatmap}
\end{figure}

\subsection{Estados internos da rede}
\label{sec:internos}

O arquivo \texttt{.raw} da variante ideal registra os 120 neurônios
ocultos, o que permite observar o mecanismo interno da rede e
relacioná-lo aos regimes de operação do amplificador operacional.

A Figura~\ref{fig:estados} mostra uma fotografia instantânea: a tensão
de saída de cada neurônio oculto, camada por camada, no instante de
medição da janela do dígito 3. Cada painel é uma camada (48, 32, 24 e
16 neurônios) e cada barra é a tensão de um neurônio; a linha
tracejada marca o teto de \SI{4,8}{\volt}. A distribuição observada é
a esperada para uma rede com ReLU: a maioria dos neurônios está em
\SI{0}{\volt} (amplificador saturado no trilho inferior, neurônio
cortado pela ativação), uma fração opera na região linear com tensões
intermediárias e poucos neurônios encostam no teto (saturação
superior). Os três grupos correspondem exatamente aos três trechos da
curva de transferência da Figura~\ref{fig:vtc}.

\begin{figure}[htb]
  \centering
  \includegraphics[width=\textwidth]{estados_internos.png}
  \caption{Tensão de saída dos 120 neurônios ocultos, camada por
  camada, no instante de medição da janela do dígito 3 (variante
  ideal). A linha tracejada é o teto de saturação de
  \SI{4,8}{\volt}.}
  \label{fig:estados}
\end{figure}

A Figura~\ref{fig:atividade} acompanha a última camada oculta (16
neurônios) ao longo das vinte janelas. Cada coluna é uma janela e cada
linha um neurônio; a intensidade é a tensão no instante de medição.
Observa-se que cada dígito ativa uma combinação própria de neurônios e
que as duas apresentações do mesmo dígito (janelas 0 e 10, 1 e 11,
etc.) produzem colunas semelhantes, embora não idênticas, pois as
imagens são amostras diferentes da mesma classe. É esse código
interno, e não qualquer pixel individual, que a camada de saída
converte em um único comparador ativo.

\begin{figure}[htb]
  \centering
  \includegraphics[width=\textwidth]{atividade_oculta.png}
  \caption{Tensão dos 16 neurônios da última camada oculta em cada uma
  das vinte janelas (variante ideal): cada dígito ativa um padrão
  próprio, repetido de forma aproximada na segunda apresentação da
  mesma classe.}
  \label{fig:atividade}
\end{figure}

A Figura~\ref{fig:terra} examina um único neurônio, $a^{(4)}_7$, sob
dois pontos de vista simultâneos. O painel superior mostra a tensão de
saída ao longo de dez janelas: o neurônio alterna entre janelas em que
opera na região linear (tensões entre \SI{0}{\volt} e
\SI{4,8}{\volt}) e janelas em que a ativação o corta em
\SI{0}{\volt}. O painel inferior mostra o nó de soma do mesmo
neurônio, em escala de milivolts. Enquanto a saída está na região
linear, a realimentação mantém o nó de soma no terra virtual, a
milivolts do zero -- condição que garante que a soma ponderada seja de
fato uma soma de correntes definidas apenas pelos resistores. Quando a
saída satura em \SI{0}{\volt}, a malha de realimentação abre
(o amplificador não pode mais baixar a saída) e o nó de soma
afasta-se do zero em algumas centenas de milivolts, seguindo as
entradas. A saturação que implementa a ativação e a perda do terra
virtual são o mesmo evento observado em dois nós distintos do
circuito; o erro de soma associado não se propaga, pois ocorre
justamente quando a saída está grampeada no valor correto da ReLU.

\begin{figure}[htb]
  \centering
  \includegraphics[width=\textwidth]{terra_virtual.png}
  \caption{Neurônio $a^{(4)}_7$ (variante ideal): tensão de saída
  (acima) e nó de soma em escala de milivolts (abaixo). Na região
  linear o nó permanece em terra virtual; quando a saída satura em
  \SI{0}{\volt}, a malha abre e o nó de soma afasta-se do zero.}
  \label{fig:terra}
\end{figure}

\subsection{Comparação com o modelo em software}
\label{sec:medidas}

A Tabela~\ref{tab:medidas} compara, para as dez primeiras janelas, a
saída $z$ do dígito correto em três avaliações: o modelo em software
(numpy, com os pesos E96), o circuito com amplificador ideal e o
circuito com LM741. A coluna ``maior errada'' registra a maior saída
entre os nove dígitos incorretos (medida na variante LM741) e a coluna
``disparo'' indica se apenas o comparador correto ultrapassou o limiar
naquela janela.

\begin{table}[htb]
  \centering
  \caption{Saída $z$ do dígito correto em cada janela: modelo em
  software, circuito com amplificador ideal e circuito com LM741.
  ``Maior errada'' é a maior saída entre os dígitos incorretos e
  ``disparo'' indica disparo único correto (ambos na variante LM741).}
  \label{tab:medidas}
  % Gerado por analyze_results.py - nao editar a mao.
\begin{tabular}{c S[table-format=1.3] S[table-format=1.3] S[table-format=1.3] S[table-format=-1.3] c}
\toprule
{Dígito} & {$z$ modelo (\si{\volt})} & {$z$ ideal (\si{\volt})} & {$z$ LM741 (\si{\volt})} & {maior errada (\si{\volt})} & {Disparo} \\
\midrule
0 & 4.800 & 4.800 & 4.817 & -2.951 & Sim \\
1 & 4.800 & 4.799 & 4.817 & -1.617 & Sim \\
2 & 4.800 & 4.799 & 4.817 & -1.020 & Sim \\
3 & 4.800 & 4.800 & 4.817 & -0.126 & Sim \\
4 & 4.800 & 4.799 & 4.817 & -2.149 & Sim \\
5 & 4.800 & 4.799 & 4.817 & 0.292 & Sim \\
6 & 4.800 & 4.800 & 4.816 & -2.187 & Sim \\
7 & 4.800 & 4.799 & 4.817 & -2.750 & Sim \\
8 & 4.800 & 4.800 & 4.817 & -2.586 & Sim \\
9 & 4.800 & 4.800 & 4.817 & -0.408 & Sim \\
\bottomrule
\end{tabular}

\end{table}

A Figura~\ref{fig:scatter} estende a comparação às 200 medidas de $z$
de cada variante (dez saídas em vinte janelas). O eixo horizontal é o
valor previsto pelo modelo e o vertical o valor medido no LTspice; a
reta cinza é a identidade, sobre a qual estariam medidas perfeitas.
Com o amplificador ideal os pontos são indistinguíveis da identidade
(erro médio de \vErrIdealMedio~\si{\milli\volt}, máximo de
\vErrIdealMax~\si{\milli\volt}): esse resíduo mede o erro estrutural
do circuito, dominado pelo arredondamento E96 e pelo ganho finito.
Com o LM741 o erro médio sobe para \vErrLmMedio~\si{\milli\volt}, mas
a dispersão concentra-se no canto inferior esquerdo do gráfico, na
região das saídas fortemente negativas dos dígitos rejeitados, onde
offset, ganho finito e a saturação real do LM741 se acumulam sem
consequência: essas saídas permanecem volts abaixo do limiar de
comparação. Nas saídas que decidem, próximas ou acima de $V_{TH}$, a
diferença para o modelo é de dezenas de milivolts, essencialmente a
folga entre o nível de saturação calibrado do LM741
($\approx\SI{4,81}{\volt}$) e o teto de \SI{4,8}{\volt} do modelo. A
tabela confirma a margem: em vinte janelas, a maior saída errada não
passou de \SI{0,29}{\volt}, mais de \SI{2,3}{\volt} abaixo do limiar.
Por essa razão a decisão é idêntica nas três avaliações:
\vDisparosLm~das vinte janelas com disparo único correto.

\begin{figure}[htb]
  \centering
  \includegraphics[width=0.66\textwidth]{logits_spice_vs_numpy.png}
  \caption{Saídas $z$ medidas no LTspice contra o modelo em software
  (200 medidas por variante). O erro do LM741 concentra-se nas saídas
  fortemente negativas, sem efeito sobre a decisão.}
  \label{fig:scatter}
\end{figure}

No conjunto de teste completo (\num{10000} imagens), avaliado em
software com os pesos E96, o modelo atinge \SI{97,15}{\percent} de
acurácia pelo critério do máximo. Pelo critério mais exigente do
comparador, \SI{96,44}{\percent} das imagens produzem disparo único
correto, \SI{0,86}{\percent} não disparam comparador algum (todas as
saídas abaixo do limiar) e \SI{0,51}{\percent} disparam mais de um. A
Figura~\ref{fig:confusao} apresenta a matriz de confusão pelo critério
do máximo: linhas são o dígito verdadeiro, colunas a resposta da rede,
e a diagonal concentra a quase totalidade das imagens. Os erros
restantes agrupam-se nos pares esperados para a resolução $8\times 8$,
como 4 confundido com 9 e 5 com 3, pares cujos traços se tornam
semelhantes após a redução da imagem.

\begin{figure}[htb]
  \centering
  \includegraphics[width=0.68\textwidth]{matriz_confusao.png}
  \caption{Matriz de confusão no conjunto de teste completo
  (\num{10000} imagens, resposta pelo neurônio de maior saída). Os
  erros concentram-se em pares visualmente próximos na resolução
  reduzida.}
  \label{fig:confusao}
\end{figure}

\subsection{Rede completa com o LT1810}
\label{sec:moderno-resultados}

As duas simulações de janela reduzida descritas na
Seção~\ref{sec:moderno-rede} verificam se as correções derivadas nas
bancadas sustentam a rede completa de 260 amplificadores. A
Figura~\ref{fig:moderno-saidas} resume o resultado no formato da
Figura~\ref{fig:heatmap}: tensão dos dez comparadores ao fim de cada
uma das seis janelas, sem compensação (acima) e com compensação
(abaixo); o contorno vermelho marca o dígito apresentado e a linha
pontilhada separa as duas janelas de partida, excluídas das
estatísticas.

O contraste é imediato. Sem compensação, o comparador do dígito 7
domina cinco das seis janelas, independentemente da imagem
apresentada; das quatro janelas válidas, apenas uma
(\vModRedeAcertosSemComp~de \vModRedeJanelasValidas) termina com o
comparador correto sozinho em nível alto, e mesmo nesse caso o
disparo ocorreu ainda na janela anterior. Com compensação, todas as
janelas, inclusive as de partida, terminam com disparo único correto
(\vModRedeAcertosComp~de \vModRedeJanelasValidas~válidas).

\begin{figure}[htb]
  \centering
  \includegraphics[width=0.85\textwidth]{moderno_rede_saidas.png}
  \caption{Rede completa com LT1810, janela reduzida de
  \SI{15}{\micro\second}: tensão dos dez comparadores ao fim de cada
  janela, sem compensação (acima) e com compensação (abaixo). O
  contorno vermelho indica o dígito apresentado; a linha pontilhada
  separa as janelas sob a rampa de partida.}
  \label{fig:moderno-saidas}
\end{figure}

A Figura~\ref{fig:moderno-ondas} mostra as formas de onda por trás
das duas matrizes. Em cada painel, o traço azul é o comparador do
dígito da janela corrente, o cinza são os demais comparadores e, no
painel superior, o traço vermelho destaca $out_7$. Sem compensação
(painel superior) o circuito exibe as duas patologias previstas pelas
bancadas. A primeira é o disparo indevido: $out_7$ sobe ainda durante
a rampa de partida e permanece dominante, porque o deslocamento de
polarização de cada estágio (\vVtcOffsetMod~\si{\volt} por estágio
unitário, Figura~\ref{fig:vtc}), acumulado ao longo dos cinco
estágios, desloca as saídas $z$ o suficiente para levar $z_7$ acima do
limiar independentemente da entrada. A segunda é a instabilidade: os
traços apresentam oscilação sustentada de alta frequência (as faixas
densas do gráfico), consequência do polo formado pela capacitância de
entrada com os \SI{100}{\kilo\ohm} de realimentação em cada um dos 260
estágios. Essa oscilação tem um custo direto de simulação: o passo de
tempo permanece na escala do período da oscilação e a simulação
consumiu \vModRedeCpuSemComp~segundos de CPU para
\SI{90}{\micro\second} simulados. No painel inferior, com as duas
correções aplicadas, as duas patologias desaparecem: não há oscilação
visível, os comparadores repousam em \SI{0}{\volt} e, a cada troca de
janela, apenas o comparador do novo dígito comuta para
\SI{4,8}{\volt}. A mesma simulação passou a consumir
\vModRedeCpuComp~segundos de CPU, cerca de 34 vezes menos: a
compensação estabiliza o circuito e, por consequência, também o
integrador numérico.

\begin{figure}[htb]
  \centering
  \includegraphics[width=\textwidth]{moderno_rede_ondas.png}
  \caption{Formas de onda dos dez comparadores nas duas simulações da
  rede completa com LT1810. Sem compensação (acima): oscilação
  sustentada e disparo indevido de $out_7$ pelo acúmulo do
  deslocamento de polarização. Com compensação (abaixo): disparos
  únicos corretos em todas as janelas.}
  \label{fig:moderno-ondas}
\end{figure}

A Figura~\ref{fig:moderno-decisao} mede a latência de decisão da rede
compensada nas quatro janelas válidas. Cada traço é o comparador do
dígito correto de uma janela, com o tempo contado a partir do início
da troca de imagem; o ponto marca o cruzamento do nível de decisão de
\SI{2,4}{\volt} (metade do nível alto). As quatro decisões ocorrem
entre \SI{1,32}{\micro\second} e \vModRedeLatMax~\si{\micro\second},
com mediana de \vModRedeLatMediana~\si{\micro\second}. O pior caso de
\vModRedeLatMax~\si{\micro\second} confirma, agora na rede completa, a
estimativa de \vLatenciaCadeiaMod~\si{\micro\second} obtida na bancada
de caminho crítico (Seção~\ref{sec:tempo}): a diferença entre a
bancada e a rede é inferior a \SI{2}{\percent}. Confirma também o
regime previsto: com a compensação de \SI{2}{\pico\farad}, a latência
é limitada pela constante de tempo $R_f C_f$ dos cinco estágios em
cascata, e não pelo slew rate do amplificador.

\begin{figure}[htb]
  \centering
  \includegraphics[width=0.85\textwidth]{moderno_rede_decisao.png}
  \caption{Latência de decisão da rede completa compensada com LT1810:
  comparador do dígito correto nas quatro janelas válidas, alinhadas
  ao início da troca de imagem. Os pontos marcam o cruzamento do nível
  de decisão de \SI{2,4}{\volt}.}
  \label{fig:moderno-decisao}
\end{figure}

Nas janelas válidas, o erro médio das saídas $z$ da rede compensada em
relação ao modelo é de \vModRedeErrMedio~\si{\milli\volt} (máximo de
\vModRedeErrMax~\si{\milli\volt}), da mesma ordem do LM741. O resíduo
deve-se à parcela das correntes de polarização que o resistor de
equilíbrio não cancela (a corrente de offset entre as duas entradas) e
ao próprio deslocamento dos comparadores, que permanecem sem
correção; como as margens de decisão são de volts, o disparo não é
afetado em nenhuma janela.

\subsection{Tempo de inferência e comparação com software}
\label{sec:tempo}

No circuito, o tempo de inferência é o tempo de acomodação: o
intervalo entre a troca da imagem na entrada e o cruzamento do nível
de decisão pelo comparador correto. A Figura~\ref{fig:transicao}
amplia uma troca de janela nas duas variantes da rede completa
principal. O painel superior mostra a saída $z$ do dígito da nova
janela e o inferior o comparador correspondente. Com o amplificador
ideal (azul), a saída acompanha a rampa de entrada de
\SI{10}{\micro\second} quase sem atraso adicional e o comparador
comuta assim que $z$ cruza o limiar. Com o LM741 (laranja), a saída
sobe em rampas retas de aproximadamente
\SI{0,5}{\volt\per\micro\second} -- o slew rate do amplificador, que
se torna o fator dominante porque cada um dos cinco estágios em
cascata precisa percorrer a sua excursão nessa taxa antes de o
seguinte concluir a dele.

\begin{figure}[htb]
  \centering
  \includegraphics[width=0.9\textwidth]{transicao_janela.png}
  \caption{Ampliação de uma troca de imagem: saída $z$ do dígito da
  nova janela (acima) e comparador correspondente (abaixo), com o
  amplificador ideal e com o LM741. As rampas retas do LM741 são o seu
  slew rate de \SI{0,5}{\volt\per\micro\second}.}
  \label{fig:transicao}
\end{figure}

Medida sobre as formas de onda das vinte janelas, a latência de
decisão do LM741 tem mediana de
\vLatenciaLmMediana~\si{\micro\second} e pior caso de
\vLatenciaLmMax~\si{\micro\second}; a variante ideal, limitada apenas
pela rampa de entrada e pelo RC de saída do macromodelo, decide com
mediana de \vLatenciaIdealMediana~\si{\micro\second} e pior caso de
\vLatenciaIdealMax~\si{\micro\second}.

A Figura~\ref{fig:cadeia} apresenta as bancadas de caminho crítico:
cinco somadores unitários em série e o comparador, submetidos a um
degrau de entrada, com o LM741 (esquerda) e com o LT1810 compensado
(direita), cada qual na sua escala de tempo. Em ambos os painéis o
traço azul é a saída do quinto estágio e o verde o comparador; a linha
tracejada é o limiar. Com o LM741 a cadeia decide em
\vLatenciaCadeiaLm~\si{\micro\second}, abaixo do pior caso de
\vLatenciaLmMax~\si{\micro\second} da rede completa: o tempo do LM741
é dominado pelo slew rate e as excursões da bancada
(\SI{3}{\volt} por estágio) são menores que as excursões de até
\SI{9}{\volt} das saídas da rede. A bancada fornece portanto a ordem
de grandeza correta e o mecanismo dominante, e o valor exato requer a
rede completa. Com o LT1810 compensado a cadeia decide em
\vLatenciaCadeiaMod~\si{\micro\second}, valor confirmado pela rede
completa compensada (\vModRedeLatMax~\si{\micro\second},
Seção~\ref{sec:moderno-resultados}).

\begin{figure}[htb]
  \centering
  \includegraphics[width=\textwidth]{cadeia_critica.png}
  \caption{Bancadas de caminho crítico (cinco somadores unitários e
  comparador) respondendo a um degrau de entrada: LM741 (esquerda) e
  LT1810 compensado (direita), em escalas de tempo distintas. A linha
  tracejada é o limiar $V_{TH}$.}
  \label{fig:cadeia}
\end{figure}

No computador usado como referência (numpy nativo no hospedeiro), o
mesmo modelo responde em \SI{\vMacUmaImagem}{\micro\second} para uma
imagem avulsa e em \SI{\vMacLote}{\micro\second} por imagem quando
processa as \num{10000} imagens em lote. A Tabela~\ref{tab:fps} e a
Figura~\ref{fig:tempos} traduzem os piores casos medidos em taxa de
imagens por segundo, adotando o pior caso como período seguro de
amostragem. Para inferência de imagem única, a rede com LM741 fica
uma ordem de grandeza abaixo do computador; a rede com LT1810
compensado, decidindo em menos de \SI{2}{\micro\second}, supera a
inferência avulsa do computador por fator próximo de cinco. Em lote,
o computador amortiza custos fixos sobre \num{10000} imagens e supera
qualquer das variantes por mais de uma ordem de grandeza.

\begin{table}[htb]
  \centering
  \caption{Tempo por imagem e taxa correspondente de cada
  implementação. Para os circuitos, o período é o pior caso medido da
  latência de decisão (rede completa em ambos os casos).}
  \label{tab:fps}
  \begin{tabular}{lrr}
    \toprule
    Implementação & Tempo por imagem & Imagens por segundo \\
    \midrule
    Rede com LM741 & \vLatenciaLmMax~\si{\micro\second} & \vFpsLm \\
    Rede com LT1810 compensado & \vModRedeLatMax~\si{\micro\second} & \vFpsModRede \\
    Mac, uma imagem por vez & \vMacUmaImagem~\si{\micro\second} & \vMacFpsUma \\
    Mac, lote de \num{10000} & \vMacLote~\si{\micro\second} & \vMacFpsLote \\
    \bottomrule
  \end{tabular}
\end{table}

\begin{figure}[htb]
  \centering
  \includegraphics[width=\textwidth]{tempo_por_imagem.png}
  \caption{Tempo por imagem de cada implementação, em escala
  logarítmica. Para o LT1810 a figura usa a estimativa da bancada de
  caminho crítico (\vLatenciaCadeiaMod~\si{\micro\second}), que difere
  menos de \SI{2}{\percent} do valor medido na rede completa
  compensada.}
  \label{fig:tempos}
\end{figure}

Quanto ao custo, a análise limita-se a ordens de grandeza. Uma
montagem discreta desta rede (260 amplificadores, cerca de 6000
resistores de precisão e as placas correspondentes) situa-se na casa
de poucos milhares de reais e ocupa área de bancada considerável. A
integração da mesma topologia em um processo maduro de
\SI{180}{\nano\metre}, com amplificadores compactos e resistores
programáveis armazenando os pesos, resultaria em um circuito integrado
de algumas dezenas de milímetros quadrados: unidades de poucos dólares
em volume, porém com custo fixo de projeto e máscaras na casa das
centenas de milhares de dólares. Em contrapartida, um microcontrolador
de baixo custo executa o mesmo modelo em software em alguns
milissegundos. Conclui-se que a implementação analógica só se
justifica quando latência de microssegundos e consumo muito baixo são
requisitos simultâneos; como instrumento didático, entretanto, ela tem
o mérito da concretude, pois cada peso da rede é um resistor
identificável no esquemático e cada não idealidade do amplificador tem
efeito mensurável e explicável.
